\section{Methods}
\label{sec:methods}

\subsection{Overview}

PowerTrace-Sim predicts the node power of an inference server from a request
trace and the server configuration.  Request $i$ is represented by its arrival
time $a_i$, prompt length $n_i^{\mathrm{in}}$, and output length
$n_i^{\mathrm{out}}$.  A timing model first replays these requests through a
continuous-batching scheduler.  The replay produces the start time, duration,
and work of each engine iteration.  A power model then maps this work to power
at 250-ms resolution.

This separation is important.  Arrival rate and model name do not enter the
power regression.  Their effects appear through the schedule and the work it
creates.  The same timing model is used for dense and mixture-of-experts (MoE)
models, while their power models are fitted separately because sparse expert
execution changes weight traffic and device occupancy.

The prediction is the sum of power over the $p$ GPUs in the tensor-parallel
group.  All filtering, interpolation, and delay operations are confined to a
single run.  The experiments use observed output lengths.  Prospective use
requires an output-length estimate or distribution.

\subsection{Measurements and Alignment}
\label{sec:power-alignment}

Each run contains four synchronized records: request measurements, engine
counters, GPU telemetry, and a manifest.  Request records provide arrival
times, prompt and output token counts, time to first token (TTFT), and either
per-token latency or mean inter-token latency.  Engine counters report running
and waiting requests, prompt and generation tokens, iteration counts, tokens
per iteration, and KV-cache use.  GPU power, clocks, utilization, memory use,
and temperature are sampled with \texttt{nvidia-smi} at approximately 4~Hz.
The manifest records the model architecture, GPU type, tensor-parallel size,
data types, scheduler limits, cache policy, software version, and active GPU
UUIDs.

Prompt length is a logical request quantity.  When prefix caching is enabled,
the executed prefill length is
\begin{equation}
  n_i^{\mathrm{exec}} = n_i^{\mathrm{in}}-n_i^{\mathrm{cache}},
\end{equation}
where $n_i^{\mathrm{cache}}$ is reported by the server.  Cached tokens remain
part of the request context but do not incur prefill computation.  Cache policy
and cache-hit counters are checked against the manifest before a run is used.

Request and engine timestamps are Unix epochs.  GPU timestamps are converted
to the same basis using the UTC offset stored in the manifest.  For legacy
traces, let $t_0^{\mathrm{req}}$ and $t_0^P$ be the first request and power
timestamps.  Their residual offset is
\begin{equation}
  \delta=(t_0^{\mathrm{req}}-t_0^P)-1800K,
  \qquad
  K=\operatorname{round}\!\left(
  \frac{t_0^{\mathrm{req}}-t_0^P}{1800}\right).
  \label{eq:clock-alignment}
\end{equation}
The half-hour correction detects timezone errors; every retained run has
$K=0$.  Runs are rejected when the remaining offset lies outside
$[-2,600]$ seconds.  Alignment is determined from timestamps, never by shifting
one power trace to resemble another.

A power sample at time $t_j^P$ is placed at
$\tau_j=t_j^P-t_0^P-\delta$.  Samples are matched to the active GPU UUIDs,
summed across the tensor-parallel group, and averaged within each 250-ms bin.
Bins without a sample remain missing.  Source-file hashes are stored with the
joined trace.

\subsection{From Requests to Engine Work}
\label{sec:timing-model}

\subsubsection{Work in one iteration}

The timing model calculates work from the architecture rather than fitting a
throughput for each checkpoint.  Consider an iteration that decodes one token
for each of $d$ sequences and processes a prefill chunk of $q$ tokens.  Let
$N_{\mathrm{act}}$ be the active parameter count, $L$ the number of layers,
$n_q$ the number of query heads, and $h$ the head dimension.  Mixed-precision
work is scaled by
\begin{equation}
  s_{\mathrm{dtype}}=1-\frac{1}{2}
  \operatorname{clip}(f_{\mathrm{FP8}},0,1),
  \label{eq:dtype-scale}
\end{equation}
where $f_{\mathrm{FP8}}$ is the fraction of floating-point work executed at
FP8 throughput.

Let $C(c)$ be the visible attention context.  For full attention,
$C(c)=\max(c,0)$.  For a model with equal numbers of full-attention and
sliding-window layers,
\begin{equation}
  C(c)=\frac{1}{2}\max(c,0)+
       \frac{1}{2}\min\{\max(c,0),w\},
  \label{eq:effective-context}
\end{equation}
where $w$ is the window length.  If $c_d$ is the mean decode context and $c_p$
is the context preceding the current prefill chunk, the linear and attention
work are
\begin{align}
  F_{\mathrm{lin}}
    &=2s_{\mathrm{dtype}}N_{\mathrm{act}}(d+q),\\
  F_{\mathrm{attn}}
    &=4s_{\mathrm{dtype}}n_qhL
      \left[dC(c_d)+qC\!\left(c_p+\frac{q}{2}\right)\right].
  \label{eq:iteration-flops}
\end{align}

For BF16 KV storage, the bytes stored per context token are
\begin{equation}
  v_{\mathrm{KV}}=2L n_{\mathrm{KV}}h\times2\ \mathrm{bytes},
  \label{eq:kv-bytes}
\end{equation}
and the attention traffic in one iteration is
\begin{equation}
  M_{\mathrm{attn}}=v_{\mathrm{KV}}
  \left[d\{C(c_d)+1\}+C\!\left(c_p+\frac{q}{2}\right)+q\right].
  \label{eq:attention-bytes}
\end{equation}
The prefill term represents one streaming pass through visible KV blocks and
the write of the new chunk.  It does not charge a separate full-cache read for
every query token within a fused attention kernel.

\subsubsection{Iteration time}

Linear and attention operators execute in sequence, so each receives its own
roofline bound.  Let $F_h$ and $B_h$ be the peak compute rate and HBM bandwidth
per GPU on hardware $h$, and let $\eta_F$ and $\eta_B$ be fitted efficiencies.
For expected weight traffic $M_{\mathrm{w}}$, iteration time is
\begin{align}
  T_{\mathrm{iter}}={}&T_{\mathrm{launch}}
  +\max\!\left(
    \frac{F_{\mathrm{lin}}}{p\eta_FF_h},
    \frac{M_{\mathrm{w}}}{p\eta_BB_h}
  \right) \nonumber\\
  &+\max\!\left(
    \frac{F_{\mathrm{attn}}}{p\eta_FF_h},
    \frac{M_{\mathrm{attn}}}{p\eta_BB_h}
  \right)+dT_{\mathrm{sample}}.
  \label{eq:iteration-time}
\end{align}
The launch and synchronization term is
\begin{equation}
  T_{\mathrm{launch}}=T_0+2L T_{\mathrm{msg}}(h,p).
  \label{eq:launch-time}
\end{equation}
The factor $2L$ represents two tensor-parallel collectives per layer; for
$p=1$, the fitted term captures the corresponding kernel-launch chain.

The efficiencies and overheads are shared by all checkpoints on one hardware
type.  They are fitted in log-time space from controlled prefill and decode
probes, isolated production requests, and loaded requests grouped by measured
decode concurrency.  Held-out model families are excluded.  FP8 memory
throughput is calibrated separately on H100.

\subsubsection{Continuous batching}

Requests enter a first-come-first-served queue.  A request is admitted when a
sequence slot and sufficient KV capacity are available.  If $g$ is the
recorded GPU-memory fraction and $M_h$ is memory capacity per GPU, the available
context capacity is
\begin{equation}
  Q_{\mathrm{KV}}=
  \frac{g pM_h-M_{\mathrm{weights}}}{v_{\mathrm{KV}}}.
  \label{eq:kv-capacity}
\end{equation}
The simulation uses the recorded sequence limit and token budget.  Each
iteration first advances all sequences ready to decode, then assigns the
remaining token budget to prefill chunks in arrival order.  The simulation
records iteration time, decode batch, context length, prefill chunks, weight
traffic, FLOPs, and attention bytes.

Iteration intervals are projected onto all overlapping 250-ms bins in
proportion to overlap duration.  This preserves total work, bytes, and busy
time across bin boundaries.  Token completions and iteration counts are placed
at iteration completion.  The evaluation horizon is fixed by the measured
request trace: an early simulated drain produces idle bins, while simulated
events beyond the common horizon are discarded.  Measured power is not used
to construct any model input.

\subsection{Weight Traffic for MoE Models}
\label{sec:routing-method}

A dense iteration reads the resident weight set.  For an MoE model, let $E$ be
the number of experts, $k$ the number selected per token, and $f_{\mathrm{MoE}}$
the fraction of weight bytes assigned to routed experts.  Under independent
uniform top-$k$ routing, the expected fraction of experts touched by
$z=d+q$ tokens is
\begin{equation}
  \phi(z)=1-\left(1-\frac{k}{E}\right)^z.
  \label{eq:uniform-routing}
\end{equation}
Expected weight traffic is therefore
\begin{equation}
  M_{\mathrm{w}}=M_{\mathrm{weights}}
  \left[(1-f_{\mathrm{MoE}})+f_{\mathrm{MoE}}\phi(z)\right].
  \label{eq:moe-weight-bytes}
\end{equation}
This law is used for the reported MoE timing and power results.  Captured
expert assignments were evaluated separately, but did not support a stable
additional routing term.

\subsection{Dense Power Model}
\label{sec:dense-power}

The dense model is fitted once per GPU type.  Let
$r_{F,\mathrm{lin}}(t)$ and $r_{F,\mathrm{attn}}(t)$ be the linear and
attention FLOP rates in bin $t$, and let $r_{\mathrm{w}}(t)$ and
$r_{B,\mathrm{attn}}(t)$ be weight and attention byte rates.  Per-GPU
compute and memory utilization are
\begin{align}
  u_F(t)&=\frac{r_{F,\mathrm{lin}}(t)+r_{F,\mathrm{attn}}(t)}{pF_h},\\
  u_B(t)&=\frac{r_{\mathrm{w}}(t)+r_{B,\mathrm{attn}}(t)}{pB_h}.
  \label{eq:dense-utilization}
\end{align}
The hardware constants are
\begin{align}
  (F_{\mathrm{A100}},B_{\mathrm{A100}})
    &=(312\ \mathrm{TFLOP/s},2.0\ \mathrm{TB/s}),\\
  (F_{\mathrm{H100}},B_{\mathrm{H100}})
    &=(990\ \mathrm{TFLOP/s},3.35\ \mathrm{TB/s}).
\end{align}

Let $b(t)\in[0,1]$ be the fraction of the bin occupied by engine work.  The
complete per-GPU design has four coordinates:
\begin{equation}
  \mathbf{x}_D(t)=\left[
  1,\;
  b(t)\operatorname{clip}\!\left(
    \frac{M_{\mathrm{weights}}}{pM_h},0,1\right),\;
  u_F(t),\;
  \sqrt{b(t)\max\{u_B(t),0\}}
  \right].
  \label{eq:dense-design}
\end{equation}
These terms represent idle power, the active resident model, compute demand,
and memory demand.  The square-root memory response is applied before active
and idle time are combined, which preserves short memory-intensive intervals.
The model contains no checkpoint name, arrival rate, tensor-parallel label,
phase label, or elapsed-time term.  Nonnegative coefficients make predicted
power monotone in each work coordinate.

For hardware $h$, node power is
\begin{equation}
  \widehat P_D(t)=p\,
  \{\mathcal H_h(\mathbf X_D)\}_t\boldsymbol\beta_{D,h},
\end{equation}
where $\mathcal H_h$ is the measured sensor response applied independently to
each run.

\subsection{MoE Power Model}
\label{sec:moe-power}

The MoE model uses the same simulated work but retains a separate regression.
Let
\begin{equation}
  m(t)=r_{\mathrm{w}}(t)+r_{\mathrm{KV,r}}(t)+r_{\mathrm{KV,w}}(t)
\end{equation}
be node-level memory traffic, $r_{\mathrm{iter}}(t)$ the iteration rate, and
$b_{\mathrm{dec}}(t)$ the mean decode batch.  The per-GPU design is
\begin{equation}
  \mathbf{x}_M(t)=\left[
    1,\;
    \mathbb{1}(p>1),\;
    \frac{m(t-\Delta t)}{p(2.0\times10^{12})},\;
    \sqrt{b(t)u_F(t)},\;
    \frac{r_{\mathrm{iter}}(t)}{1000},\;
    \log\{1+b_{\mathrm{dec}}(t)\}
  \right].
  \label{eq:moe-design}
\end{equation}
The 250-ms memory lag represents the observed response delay; it is reset at
every run boundary.  The remaining terms represent idle power, the multi-GPU
floor, compute demand, launch activity, and the sublinear effect of decode
concurrency.

The regression form is shared, but GPT-OSS-20B and GPT-OSS-120B have separate
coefficients because their observed tensor-parallel ranges do not overlap.
The MoE fit is therefore architecture-specific rather than an extrapolation
across unsupported tensor-parallel configurations.

\subsection{Calibration}
\label{sec:power-fit}

Both power models use nonnegative least squares.  The dense idle coefficient
is measured before fitting.  Within each source run, zero-busy intervals of at
least four seconds are located, their first two seconds are discarded, and
the remaining median is recorded.  The hardware idle floor is the median over
runs.  This prevents long traces from dominating the baseline.

Dense fitting uses matched, nonoverlapping one-second means.  For each run $r$,
let $q_r$ be its 90th-percentile power, and let $\mathcal V_r^-$ and
$\mathcal V_r^+$ contain the seconds below and at or above $q_r$.  With
$n_r^s=|\mathcal V_r^s|$, the fit is
\begin{equation}
  \widehat{\boldsymbol\beta}_{D,h}=
  \mathop{\arg\min}_{\substack{\boldsymbol\beta\ge0\\
    \beta_0=P_{\mathrm{idle},h}}}
  \sum_{r\in\mathcal R_{D,h}}
  \sum_{s\in\{-,+\}}\frac{1}{2n_r^s}
  \sum_{t\in\mathcal V_r^s}
  \left[P_r(t)-
  \{\mathcal H_h(\mathbf X_r)\}_t\boldsymbol\beta\right]^2.
  \label{eq:dense-objective}
\end{equation}
Every run has equal weight, divided equally between ordinary and upper-tail
power.  Brief prefill peaks therefore remain visible beside long decode
periods.  The response delay is selected from
$\{0,0.25,0.5,0.75\}$ seconds with the same objective.  A100 selects no delay;
H100 selects 0.25 seconds and a one-second trailing mean.

The MoE fit gives equal total weight to every source run at 250-ms resolution:
\begin{equation}
  \widehat{\boldsymbol\beta}_{M}=
  \arg\min_{\boldsymbol\beta\ge0}
  \sum_{r\in\mathcal R_M}\frac{1}{n_r}
  \sum_{t\in\mathcal V_r}
  \left[P_r(t)-\mathbf{x}_{M,r}(t)^\mathsf T
  \boldsymbol\beta\right]^2.
  \label{eq:moe-objective}
\end{equation}
Predictions are not clipped to a nominal board-power rating.  A power cap is
valid only when the corresponding operating limit is recorded for the run.

The calibration data are available from ordinary serving infrastructure:
request logs, vLLM metrics, and GPU telemetry from \texttt{nvidia-smi} or
DCGM.  Complete telemetry is needed for calibration runs, not for every
production request.  A new deployment requires controlled prefill and decode
probes, a settled loaded-idle interval, and mixed-load runs that span its
operating range.

\subsection{Data Separation}
\label{sec:data-splits}

The timing corpus contains 450 runs.  Repetitions within a
model--hardware--tensor-parallel--rate cell replay the same ordered request
lengths and nearly identical arrivals.  Splitting these repetitions between
fitting and validation would leak workload identity.  All non-twin runs at
rates up to two requests/s therefore form one source population: 225 dense
runs and 60 MoE runs.  The 57 runs at four requests/s are retrospective stress
tests.  The 108 related-model twins are transfer comparisons, not independent
request-trace validation.

Dense and MoE fitting populations are disjoint, and no rate-four target enters
either fit.  New traces from other checkpoints and arrival processes are used
only for transfer evaluation.  Each fitted artifact records the training run
identifiers and hashes of the joined ledger and timing calibration.

\subsection{Evaluation}
\label{sec:power-metrics}

Missing power samples are interpolated only for evaluation, and their count is
reported.  Measured and predicted traces are reduced to matched,
nonoverlapping one-second means.  Temporal metrics require at least 62 seconds,
at least 95\% observed one-second bins before interpolation, and no power gap
longer than two seconds.  A run that fails this gate retains its energy result
but receives no temporal score.

Let $y_t$ and $\widehat y_t$ be the measured and predicted one-second node
power for a run of length $T$.  Relative energy error is
\begin{equation}
  E_{\mathrm{err}}=100
  \frac{\left|\sum_{t=1}^T\widehat y_t-\sum_{t=1}^Ty_t\right|}
       {\sum_{t=1}^Ty_t}.
  \label{eq:energy-error}
\end{equation}
We report RMSE in watts per GPU and normalize it by the measured power range:
\begin{equation}
  \operatorname{NRMSE}_{\mathrm{range}}=
  \frac{\sqrt{T^{-1}\sum_{t=1}^T(\widehat y_t-y_t)^2}}
       {\max_t y_t-\min_t y_t}.
  \label{eq:nrmse}
\end{equation}

For $z_t=y_t-\overline y$, the autocorrelation at lag $\ell$ is
\begin{equation}
  \rho_y(\ell)=
  \frac{\sum_{t=1}^{T-\ell}z_tz_{t+\ell}}
       {\sum_{t=1}^{T}z_t^2},
  \qquad \ell=1,\ldots,60.
  \label{eq:acf}
\end{equation}
The predicted autocorrelation is defined identically.  We report its mean
absolute error and
\begin{equation}
  R^2_{\mathrm{ACF}}=1-
  \frac{\sum_{\ell=1}^{60}
  [\rho_{\widehat y}(\ell)-\rho_y(\ell)]^2}
  {\sum_{\ell=1}^{60}
  [\rho_y(\ell)-\overline{\rho_y}]^2}.
  \label{eq:acf-r2}
\end{equation}
Distributional agreement is $1-D_{\mathrm{KS}}$, where $D_{\mathrm{KS}}$ is
the two-sample Kolmogorov--Smirnov distance between measured and predicted
per-GPU power.

Soft-DTW is reported as a temporal diagnostic.  Both traces are centered by
the measured mean and scaled by the measured range.  With smoothing
$\gamma=0.01$ and a ten-second Sakoe--Chiba band $B$, the reported divergence is
\begin{equation}
  D_{\gamma,B}(u,v)=\frac{1}{T}\left[
  \operatorname{sDTW}_{\gamma,B}(u,v)
  -\frac{1}{2}\operatorname{sDTW}_{\gamma,B}(u,u)
  -\frac{1}{2}\operatorname{sDTW}_{\gamma,B}(v,v)\right].
  \label{eq:soft-dtw}
\end{equation}
Soft-DTW is not used to fit coefficients or select a model.  All metrics are
computed per run and summarized across runs; time bins from different runs are
never pooled.

\subsection{Reproducibility}

The request ledger, power join, and model fit are produced by scripts under
version control.  Stored coefficients are evaluated without refitting.
Run identifiers, split labels, source hashes, and the timing artifact hash are
stored with each power artifact.  The complete commands are provided with the
artifact.
